\documentclass{article}
\usepackage{mathtools} 
\usepackage{fontspec}
\usepackage[UTF8]{ctex}
\usepackage{amsthm}
\usepackage{mdframed}
\usepackage{xcolor}
\usepackage{amssymb}
\usepackage{amsmath}


% 定义新的带灰色背景的说明环境 zremark
\newmdtheoremenv[
  backgroundcolor=gray!10,
  % 边框与背景一致，边框线会消失
  linecolor=gray!10
]{zremark}{说明}

% 通用矩阵命令: \flexmatrix{矩阵名}{元素符号}{行数}{列数}
\newcommand{\flexmatrix}[4]{
  \[
  #1 = \begin{pmatrix}
    #2_{11}     & #2_{12}     & \cdots & #2_{1#4}   \\
    #2_{21}     & #2_{22}     & \cdots & #2_{2#4}   \\
    \vdots      & \vdots      & \ddots & \vdots     \\
    #2_{#31}    & #2_{#32}    & \cdots & #2_{#3#4}
  \end{pmatrix}
  \]
}

% 简化版命令（默认矩阵名为A，元素符号为a）: \quickmatrix{行数}{列数}
\newcommand{\quickmatrix}[2]{\flexmatrix{A}{a}{#1}{#2}}

\begin{document}
\title{注释 18.3}
\author{张志聪}
\maketitle
\begin{zremark}
  空间直线的表示方式
\end{zremark}

\begin{itemize}
  \item 一般式，即两平面相交的直线；

        \begin{equation*}
          \begin{cases}
            a x + b y + c z = 0 \\
            d x + e y + f z = 0
          \end{cases}
        \end{equation*}
  \item 点向式；

        已知直线上一点$P_0(x_0, y_0, z_0)$，直线方向
        $\vec{v} = (a, b, c)$，则
        \begin{align*}
          \frac{x - x_0}{a} & = \frac{y - y_0}{b} = \frac{z - z_0}{c}
        \end{align*}

  \item 参数方程.

        \begin{equation*}
          \begin{cases}
            x = a + mt \\
            y = b + nt \\
            y = c + kt
          \end{cases}
        \end{equation*}
\end{itemize}

\begin{zremark}
  空间直线与空间曲线两者的参数方程区别
\end{zremark}
空间直线的参数方程是一次函数，
空间曲线的参数方程不是一次函数。

\begin{zremark}
  求空间曲线的切线的思想
\end{zremark}

切线是动点$Q$与$P_0$的割线极限（$Q \to P_0$），即如下极限为切线方向：
\begin{align*}
   & \lim\limits_{\Delta t \to 0} \frac{\vec{r}(t_0 + \Delta t) - \vec{r}(t_0)}{\Delta t} \\
   & = \left(
  \lim\limits_{\Delta t \to 0} \frac{x(t_0 + \Delta t) - x(t_0)}{\Delta t} ,
  \lim\limits_{\Delta t \to 0} \frac{y(t_0 + \Delta t) - y(t_0)}{\Delta t},
  \lim\limits_{\Delta t \to 0} \frac{z(t_0 + \Delta t) - z(t_0)}{\Delta t}
  \right)
\end{align*}

\end{document}